% ============================================================================
%  E — file phần: Luyện hằng ngày
%  Ngày tạo: 2026-09-07 | Sửa gần nhất: 2026-09-07 | Ôn gần nhất: chưa
% ============================================================================
\documentclass[../english.tex]{subfiles}
\begin{document}
\part{LUYỆN HẰNG NGÀY}
\begin{tomtat}
Bốn phần đầu cho hiểu biết; phần này biến hiểu biết thành thói quen. Ba chương: cân bằng bốn nhánh luyện tập để không dồn hết vào một mặt (19); chọn và dùng công cụ, kể cả mô hình ngôn ngữ, mà không bị chúng làm thay (20); và một lộ trình mười hai tuần cùng ngân hàng tài nguyên (21). Nếu chỉ nhớ một điều: khoảng cách giữa \emph{biết} và \emph{dùng được} chỉ đóng lại bằng số lần bạn thật sự dùng, nên hãy thiết kế cuộc sống sao cho tiếng Anh nằm trên đường đi của bạn chứ không nằm ngoài nó.
\end{tomtat}

Có một hiện tượng quen thuộc mà bạn có lẽ đã gặp: người học đọc rất nhiều sách về tiếng Anh, biết rành các quy tắc, và vẫn viết chậm và nói ngập ngừng. Đây không phải nghịch lý. Hiểu biết về ngôn ngữ và khả năng dùng ngôn ngữ là hai thứ được lưu trữ và truy xuất theo hai cách khác nhau, và cái thứ hai chỉ hình thành qua sử dụng lặp lại.

Vì vậy phần này khác ba phần trước về tính chất. Nó không thêm nội dung mới về ngôn ngữ mà nói về \emph{thiết kế}: thiết kế một tuần luyện tập, thiết kế bộ công cụ, thiết kế một lộ trình có thể duy trì được khi bạn bận.

Chương 19 đưa ra khung cân bằng quan trọng nhất --- bốn nhánh luyện tập, với quan sát rằng phần lớn người tự học dồn gần hết thời gian vào một nhánh và bỏ trống ba nhánh còn lại. Chương 20 nói về công cụ, và dành phần lớn cho câu hỏi mà mọi người học hiện nay đều phải trả lời: dùng mô hình ngôn ngữ thế nào để nó giúp bạn giỏi lên chứ không thay bạn làm việc. Chương 21 gộp mọi thứ thành một lộ trình mười hai tuần và một danh mục tài nguyên đã chọn lọc.
\subfile{00-map}
\subfile{19-bon-nhanh-luyen/bon-nhanh-luyen}
\subfile{20-cong-cu/cong-cu}
\subfile{21-lo-trinh-va-tai-nguyen/lo-trinh-va-tai-nguyen}
\ifSubfilesClassLoaded{\printbibliography[title={Nguồn}]}{}
\end{document}
